\documentclass[12pt]{article}
\usepackage[T1]{fontenc}
\usepackage[utf8]{inputenc}
\usepackage{lmodern}
\usepackage{textcomp}
\usepackage{enumitem}
\usepackage{geometry}
\usepackage{graphicx}
\usepackage{amsmath, amssymb}
\geometry{margin=1in}
\begin{document}
\begin{center}
Chemistry - Graduate Level Assessment 20 \
Total Marks: 40 \
Answer all questions.
\end{center}

\section*{Multiple Choice (10 marks)}
\begin{enumerate}[label=\arabic*.]
\item What is the maximum number of phases that can be at equilibrium with each other in a three component mixture?
\begin{enumerate}[label=(\alph*)]
    \item 6
    \item 9
    \item 2
    \item 4
    \item 10
\end{enumerate}
\item Calculate the de Broglie wavelength for (a) an electron with a kinetic energy of $100 \mathrm{eV}$
\begin{enumerate}[label=(\alph*)]
    \item 0.175 nm
    \item 0.080 nm
    \item 0.098 nm
    \item 0.123 nm
    \item 0.135 nm
\end{enumerate}
\item One cubic millimeter of oil is spread on the surface of water so that the oil film has an area of 1 square meter. What is the thickness of the oil film in angstrom units?
\begin{enumerate}[label=(\alph*)]
    \item 200 \ensuremath{\AA}
    \item 1 \ensuremath{\AA}
    \item 0.1 \ensuremath{\AA}
    \item 50 \ensuremath{\AA}
    \item 500 \ensuremath{\AA}
\end{enumerate}
\item Which of the following molecules is a strong electrolyte when dissolved in water?
\begin{enumerate}[label=(\alph*)]
    \item CH 3 COOH
    \item CO 2
    \item O$_{2}$
    \item H$_{2}$O
    \item NH 3
\end{enumerate}
\item Which of the following is an n-type semiconductor?
\begin{enumerate}[label=(\alph*)]
    \item Silicon
    \item Boron-doped silicon
    \item Pure germanium
    \item Arsenic-doped silicon
    \item Indium Phosphide
\end{enumerate}
\end{enumerate}

\section*{True / False (10 marks)}
\textit{Answer True or False}
\begin{enumerate}[label=\arabic*.]
\setcounter{enumi}{5}
\item True or False: The nuclide 17 O has an NMR frequency of 115.5 MHz in a 20.0 T magnetic field.
\item True or False: The temperature change of 1 mole of a diatomic gas is 65 K when 65 J of heat is added and 210 J of work is done.
\item True or False: The expected value of the number of heads when two coins are tossed is 2.5.
\item True or False: The valence of element X, which has an atomic weight of 58.7 and displaces hydrogen, is -2.
\item True or False: The carbon configuration $1 s^2 2 s^2 2 p^2$ corresponds to 20 distinct quantum states.
\end{enumerate}

\section*{Long Form Response (20 marks)}
\textit{Answer each question with a clear and complete explanation.}
\begin{enumerate}[label=\arabic*.]
\setcounter{enumi}{10}
\item Discuss the kinetic energy of a 1.0 g particle released from rest near the Earth's surface under an acceleration of 8.91 m/$s^2$ after 1.0 s, ignoring air resistance.
\item Discuss the total binding energy of carbon-12 (\_6 $C^12$) and analyze its average binding energy per nucleon, exploring the implications for nuclear stability and structure.
\end{enumerate}

\vfill
\noindent\textit{End of Paper}
\end{document}